\section{Introduction}
Quantization is now the dominant technique for deploying large language
models on commodity hardware. Yet whilst memory and latency benefits are
well-understood, the question of \emph{which capabilities degrade first}
under aggressive precision reduction remains open. Mathematical reasoning
is a natural stress test: arithmetic and symbolic problems require many
sequential non-linear approximations and amplify errors quickly.

\paragraph{Question.} How does the residual error in non-linear function
approximation by a small open-weight language model scale with the
bit-width of its parameters, and which submodules---attention or
feed-forward---contribute most to that error?

\paragraph{Contributions.} (i) A dependency-free per-channel symmetric
quantizer covering FP32, FP16, INT8, INT4, and ternary (1.58-bit) regimes.
(ii) A reproducible evaluation pipeline combining a synthetic arithmetic
benchmark with optional GSM8K loading. (iii) A submodule-sensitivity test
isolating attention from feed-forward contributions to logit error. (iv) An
end-to-end paper whose numbers and figures are regenerated from raw results
on each run; the same code rescales to TinyLlama-1.1B and Qwen2.5-Math-7B
via the GPU configuration.
